% Generado por analisis/dispersion.py -- NO EDITAR A MANO.
% Panel: 19 dia(s), ultimo 2026-10-07. Criterio: EAN de fabricante (12+ d\'igitos, sin prefijo 2).

\begin{table}[htbp]
\caption{Precio id\'entico entre pares de cadenas para el mismo EAN, todos los d\'ias del panel (19 d\'ias, \'ultimo 2026-10-07)}
\label{tab:identidad}
\centering\scriptsize\setlength{\tabcolsep}{3pt}
\begin{tabular}{@{}llrrr@{}}
\toprule
Par & Relaci\'on & EAN-d\'ias & Id\'entico & Brecha med. \\
\midrule
Metro + Wong & mismo grupo & 79\,419 & 65.5\,\% & 0.0\,\% \\
Wong + Vivanda & competidores & 36\,635 & 47.6\,\% & 1.4\,\% \\
Plaza Vea + Tottus & competidores & 59\,400 & 47.5\,\% & 0.8\,\% \\
Metro + Vivanda & competidores & 34\,551 & 34.1\,\% & 3.6\,\% \\
Metro + Plaza Vea & competidores & 42\,918 & 26.9\,\% & 5.4\,\% \\
Wong + Plaza Vea & competidores & 43\,761 & 24.0\,\% & 6.2\,\% \\
Plaza Vea + Vivanda & mismo grupo & 75\,105 & 22.2\,\% & 3.5\,\% \\
Metro + Tottus & competidores & 62\,162 & 18.5\,\% & 6.6\,\% \\
Wong + Tottus & competidores & 64\,462 & 15.0\,\% & 7.3\,\% \\
Vivanda + Tottus & competidores & 47\,162 & 13.6\,\% & 6.5\,\% \\
\bottomrule\end{tabular}\end{table}

\begin{table*}[t]
\caption{Distribuci\'on de la brecha m\'ax/m\'in del precio por EAN-d\'ia y masa en cero, todos los d\'ias del panel}
\label{tab:brecha}
\centering\scriptsize\setlength{\tabcolsep}{4pt}
\begin{tabular}{@{}lrrrrrrrrr@{}}
\toprule
Subconjunto & EAN-d\'ias & Masa en 0 & p10 & p25 & p50 & p75 & p90 & SD log \\
\midrule
$\geq 2$ cadenas & 171\,125 & 24.6\,\% & 0.0\,\% & 0.4\,\% & 6.1\,\% & 16.7\,\% & 30.9\,\% & 0.065 \\
$\geq 3$ cadenas & 88\,887 & 11.1\,\% & 0.0\,\% & 3.8\,\% & 11.1\,\% & 22.2\,\% & 35.9\,\% & 0.075 \\
las 5 cadenas & 16\,910 & 5.3\,\% & 2.5\,\% & 6.5\,\% & 14.0\,\% & 25.6\,\% & 41.0\,\% & 0.076 \\
\bottomrule\end{tabular}
\par\smallskip\raggedright\scriptsize SD log: desviaci\'on est\'andar media del log-precio dentro del EAN-d\'ia; referencia de Gorodnichenko y Talavera (2017): 0.13--0.16.
\end{table*}

\begin{table}[htbp]
\caption{Descomposici\'on de la varianza de $\log$(precio) por efectos fijos anidados, todos los d\'ias del panel}
\label{tab:varianza}
\centering\scriptsize\setlength{\tabcolsep}{3pt}
\begin{tabular}{@{}llrrr@{}}
\toprule
Orden & Nivel & $R^2$ acum. & $R^2$ incr. & \% var. intra-prod. \\
\midrule
A & dia & 0.10\,\% & 0.10\,\% & -- \\
 & producto & 98.62\,\% & 98.52\,\% & -- \\
 & categoria & 98.62\,\% & 0.00\,\% & -- \\
 & grupo & 98.63\,\% & 0.01\,\% & 0.9\,\% \\
 & cadena & 98.68\,\% & 0.05\,\% & 3.4\,\% \\
 & cadena $\times$ categoria & 98.70\,\% & 0.02\,\% & 1.4\,\% \\
\midrule
B & dia & 0.10\,\% & 0.10\,\% & -- \\
 & categoria & 38.86\,\% & 38.76\,\% & -- \\
 & grupo & 38.92\,\% & 0.06\,\% & -- \\
 & cadena & 39.16\,\% & 0.24\,\% & -- \\
 & cadena $\times$ categoria & 39.81\,\% & 0.65\,\% & -- \\
 & producto & 98.70\,\% & 58.89\,\% & -- \\
\bottomrule\end{tabular}
\par\smallskip\raggedright\scriptsize Orden A: producto (EAN$\times$d\'ia) primero; la categor\'ia queda absorbida. Orden B: sin producto hasta el final. $n$ = 491\,059 observaciones EAN$\times$cadena$\times$d\'ia.
\end{table}
